% 第11章第3节：正交对角化（实对称矩阵的谱定理、谱分解的计算与核验、特征子空间与谱投影）。

\section{正交对角化}
\label{sec:ls-spectral-theory}

上一节用特征向量基完成对角化，只要求换基矩阵$\bm V$可逆，不要求它的各列互相垂直。这样的特征基一般是斜基：坐标的平方和不再等于向量长度的平方，坐标轴之间的夹角也不再是直角。本节继续同一个分解问题，但要求换基矩阵是正交矩阵，依次回答三个问题：哪些实矩阵能在标准正交基下对角化；怎样求出并检验这种分解；特征值重复时，分解中的哪些部分仍由矩阵唯一确定。

第一个问题的答案是实对称矩阵，这就是谱定理。第二个问题给出计算路线和逐项核验的方法。第三个问题的答案是特征子空间和到它们的正交投影：特征基可以在重根空间内任意旋转，这些投影却不变；由它们还能直接读出对称矩阵的定号性。对称矩阵之外的一般方阵留给下一节的Schur形；从第~\ref{sec:ls-quadratic-extrema}节起，本节的分解将分别用于长度约束下的二次型极值、最佳低秩逼近与伸缩分离三个问题。

\subsection{实对称矩阵的谱定理}
\label{sec:ls-symmetric-spectral-decomposition}

特征方向互相垂直时，换基最简单。若$\bm q_1,\ldots,\bm q_d$是$\mathbb R^d$的标准正交基，把它们排成正交矩阵$\bm Q$，则$\bm Q^{-1}=\bm Q^\top$。于是换基不必解方程：向量$\bm x$的新坐标就是内积$c_i=\bm q_i^\top\bm x$，长度也保持不变，$\lVert\bm x\rVert_2^2=\sum_ic_i^2$（第~\ref{sec:ls-orthonormal-bases}节）。所以，若$\bm S\in\mathbb R^{d\times d}$能写成
\[
 \bm S=\bm Q\bm\Lambda\bm Q^\top,
\]
其中$\bm Q$为正交矩阵、$\bm\Lambda$为实对角矩阵，它的作用就最容易看清：先用内积读出各方向的坐标，再把第$i$个坐标乘以$\lambda_i$，最后合回去。

哪些矩阵能这样写？先看必要条件。对上式取转置，由$\bm\Lambda^\top=\bm\Lambda$得
\[
 \bm S^\top=(\bm Q^\top)^\top\bm\Lambda^\top\bm Q^\top=\bm Q\bm\Lambda\bm Q^\top=\bm S,
\]
所以能用正交矩阵对角化的实矩阵必须对称。对称矩阵也就是第~\ref{sec:ls-adjoints}节的自伴映射在标准正交坐标下的矩阵：$\langle\bm S\bm x,\bm y\rangle=\langle\bm x,\bm S\bm y\rangle$对一切$\bm x,\bm y$成立。接下来的问题是：对称性是否也是充分的？本小节证明答案是肯定的。

证明的思路是：先找到一个实特征向量，再在与它垂直的子空间中重复这一步。第一步并不平凡。示例~\ref{ex:ls-rotation-no-real-eigenvector}的旋转矩阵没有实特征向量，可见一般实矩阵的特征多项式可能没有实根；对称性的作用正是排除这种情形。

所用的工具是复数。由代数基本定理，$d$次实系数多项式$p_{\bm S}$在复数范围内至少有一个根$\lambda$；第~\ref{sec:ls-general-square-spectrum}节将系统地把这些复根当作特征值，这里只用到其中一个。行列式的排列展开与消元只用到加、减、乘、除，对复数元素同样成立，所以$\det(\lambda\bm I-\bm S)=0$时，齐次方程$(\lambda\bm I-\bm S)\bm z=\bm0$有非零的复向量解$\bm z\in\mathbb C^d$，即$\bm S\bm z=\lambda\bm z$。为了在复向量上度量长度，记$\bar{\bm z}$为逐个坐标取共轭（第~\ref{sec:ru-complex-arithmetic}节）所得的向量，则
\[
 \bar{\bm z}^\top\bm z=\sum_{i=1}^d\bar z_iz_i=\sum_{i=1}^d|z_i|^2,
\]
它是实数，且$\bm z\neq\bm0$时为正。这个量称为$\bm z$与自身的共轭内积；下面的证明靠它说明对称矩阵的$\lambda$必为实数。

确认$\lambda$为实数并取得实特征向量之后，再在正交补中重复，逐个找出互相正交的特征方向，就得到下面的结论，称为\term{谱定理}（spectral theorem）。

\begin{theorem}[谱定理]
\label{thm:ls-real-symmetric-spectral}
\label{thm:matrix-analysis-symmetric-spectral-theorem}
\label{sec:ch06-triangular-spectrum}
设$\bm S\in\mathbb R^{d\times d}$满足$\bm S=\bm S^\top$，则存在正交矩阵$\bm Q=[\bm q_1,\ldots,\bm q_d]$与实对角矩阵$\bm\Lambda=\operatorname{diag}(\lambda_1,\dots,\lambda_d)$，使
\begin{equation}
 \bm S=\bm Q\bm\Lambda\bm Q^\top,\qquad \bm S\bm q_i=\lambda_i\bm q_i\quad(i=1,\ldots,d).
 \label{eq:matrix-analysis-symmetric-eigendecomposition}
\end{equation}
对任意$\bm x=\sum_i c_i\bm q_i$，二次型展开为
\begin{equation}
 \bm x^\top\bm S\bm x=\sum_{i=1}^d\lambda_i c_i^2.
 \label{eq:matrix-analysis-quadratic-spectral-form}
\end{equation}
\end{theorem}

\begin{proof}
对维数$d$归纳。$d=1$时$\bm S=[s_{1,1}]$，取$\bm Q=[1]$、$\lambda_1=s_{1,1}$即可。设结论对$d-1$阶实对称矩阵成立，下面分三步处理$d$阶的$\bm S$。

第一步，特征值是实数。由定理前的讨论，存在复数$\lambda$与非零复向量$\bm z$使$\bm S\bm z=\lambda\bm z$。左乘$\bar{\bm z}^\top$，得
\[
 \bar{\bm z}^\top\bm S\bm z=\lambda\,\bar{\bm z}^\top\bm z .
\]
另一方面，$\bm S$的元素是实数，对$\bm S\bm z=\lambda\bm z$逐坐标取共轭得$\bm S\bar{\bm z}=\bar\lambda\bar{\bm z}$；再由$\bm S^\top=\bm S$，
\[
 \bar{\bm z}^\top\bm S\bm z=(\bm S^\top\bar{\bm z})^\top\bm z=(\bm S\bar{\bm z})^\top\bm z=\bar\lambda\,\bar{\bm z}^\top\bm z .
\]
两式相减得$(\lambda-\bar\lambda)\,\bar{\bm z}^\top\bm z=0$。由于$\bar{\bm z}^\top\bm z=\sum_i|z_i|^2>0$，只能$\lambda=\bar\lambda$，即$\lambda$是实数。

第二步，取得实特征向量。把$\bm z$按实部、虚部写成$\bm z=\bm x+\mathrm i\bm y$，其中$\bm x,\bm y\in\mathbb R^d$。$\bm S$与$\lambda$都是实的，比较$\bm S\bm x+\mathrm i\,\bm S\bm y=\lambda\bm x+\mathrm i\,\lambda\bm y$两边的实部与虚部，得$\bm S\bm x=\lambda\bm x$、$\bm S\bm y=\lambda\bm y$。$\bm z\neq\bm0$说明$\bm x$与$\bm y$不全为零，取其中非零的一个并单位化，记为$\bm q_1$，记$\lambda_1=\lambda$，则$\bm S\bm q_1=\lambda_1\bm q_1$。

第三步，正交补不变并归纳。对$\bm h\perp\bm q_1$，由对称性，
\[
 \bm q_1^\top(\bm S\bm h)=(\bm S\bm q_1)^\top\bm h=\lambda_1\bm q_1^\top\bm h=0 ,
\]
所以$\bm S$把$d-1$维子空间$\bm q_1^\perp$映到自身。取$\bm q_1^\perp$的一组标准正交基排成$d\times(d-1)$矩阵$\bm U$，则$\bm S\bm U$的各列仍在$\bm q_1^\perp$中，可以写成$\bm S\bm U=\bm U\bm S'$，其中$\bm S'=\bm U^\top\bm S\bm U$是$d-1$阶实对称矩阵。由归纳假设，$\bm S'=\bm Q'\bm\Lambda'\bm Q'^\top$，$\bm Q'$为正交矩阵。令$[\bm q_2,\ldots,\bm q_d]=\bm U\bm Q'$，则这些列两两正交、长度为$1$，都与$\bm q_1$正交，且$\bm S\bm U\bm Q'=\bm U\bm S'\bm Q'=\bm U\bm Q'\bm\Lambda'$，即$\bm S\bm q_i=\lambda_i\bm q_i$（$i=2,\ldots,d$），其中$\lambda_2,\ldots,\lambda_d$是$\bm\Lambda'$的对角元。于是$\bm Q=[\bm q_1,\ldots,\bm q_d]$是正交矩阵，$\bm S\bm Q=\bm Q\bm\Lambda$，两边右乘$\bm Q^\top$即得谱分解。

最后，把$\bm x=\sum_ic_i\bm q_i$代入二次型：交叉项$c_ic_j\bm q_i^\top\bm S\bm q_j=c_ic_j\lambda_j\bm q_i^\top\bm q_j$在$i\neq j$时为零，只剩$\sum_i\lambda_ic_i^2$。
\end{proof}

证明用到的是代数基本定理与对称性：前者保证复特征向量存在，后者先使特征值为实数，再使正交补在作用下不变。复数只在第一步出现，所得的$\bm Q$与$\bm\Lambda$都是实的。证明也没有用到二次型在单位球面上的极值；第~\ref{sec:ls-quadratic-extrema}节反过来由谱分解读出这些极值。同时交换$\bm\Lambda$的对角元与$\bm Q$的对应列，谱分解仍然成立，所以下文总把特征值排成$\lambda_1\geq\lambda_2\geq\cdots\geq\lambda_d$。

对称性还直接给出不同特征方向之间的关系。设$\bm S\bm u=\lambda\bm u$、$\bm S\bm w=\mu\bm w$，且$\lambda\neq\mu$。由$\bm S^\top=\bm S$，
\[
 \lambda\,\bm u^\top\bm w=(\bm S\bm u)^\top\bm w=\bm u^\top(\bm S\bm w)=\mu\,\bm u^\top\bm w ,
\]
于是$(\lambda-\mu)\bm u^\top\bm w=0$，只能$\bm u^\top\bm w=0$。所以实对称矩阵属于不同特征值的特征空间互相正交。这一点不依赖谱定理中$\bm Q$的具体选法，下一小节计算谱分解时会用到它。

\subsection{谱分解的计算与核验}
\label{sec:ls-spectral-decomposition-computation}

谱定理保证分解存在，但证明是逐步归纳的，没有直接给出计算步骤。对手算规模的对称矩阵$\bm S$，可以按以下次序求出$\bm Q$与$\bm\Lambda$。先解特征多项式$p_{\bm S}(z)=0$，得到全部特征值；由谱定理，它们都是实数，按重数计共有$d$个。再对每个不同的特征值$\lambda$解齐次方程$(\bm S-\lambda\bm I)\bm x=\bm0$，得到特征空间$\ker(\bm S-\lambda\bm I)$的一组基。然后在每个特征空间内部做Gram--Schmidt正交化（定理~\ref{thm:ls-gram-schmidt}）并单位化。最后把所有这样得到的向量排成$\bm Q$的列，把对应的特征值按同样的顺序放到$\bm\Lambda$的对角线上。

正交化只需在各特征空间内部进行。上一小节已经证明，不同特征值的特征空间互相正交，所以来自不同特征空间的向量自动垂直。重根空间则不同：解齐次方程得到的一组基通常不正交，必须在空间内部重新选择标准正交基。各特征空间的维数之和恰为$d$，这同样由谱定理保证：$\bm Q$的$d$列按特征值分组，每组都落在相应的特征空间中。

算出$\bm Q$与$\bm\Lambda$之后，逐项核验两件事。其一是$\bm Q^\top\bm Q=\bm I$，即各列单位长度且两两正交；它保证$\bm Q^{-1}=\bm Q^\top$。其二是$\bm S\bm Q=\bm Q\bm\Lambda$，逐列看就是$\bm S\bm q_i=\lambda_i\bm q_i$；这一步同时检查了特征值与列向量的顺序是否一致。两项都成立时，右乘$\bm Q^\top$即得$\bm S=\bm Q\bm\Lambda\bm Q^\top$。下面用一个二阶矩阵完成全过程。

\begin{example}[二阶对称矩阵的谱分解]
\label{ex:ls-symmetric-spectral-two-by-two}
给定
\[
 \bm S=\begin{bmatrix}2&1\\1&2\end{bmatrix},
\]
求它的谱分解并核验。

特征多项式为$p_{\bm S}(z)=(z-2)^2-1=(z-3)(z-1)$，特征值为$3$与$1$。解$(\bm S-3\bm I)\bm x=\bm0$得特征空间由$(1,1)^\top$张成，解$(\bm S-\bm I)\bm x=\bm0$得特征空间由$(1,-1)^\top$张成。两个特征值都是单根，各特征空间是一条直线，只需单位化：$\bm q_1=(1,1)^\top/\sqrt2$、$\bm q_2=(1,-1)^\top/\sqrt2$。于是
\[
 \bm Q=\frac1{\sqrt2}\begin{bmatrix}1&1\\1&-1\end{bmatrix},\qquad
 \bm\Lambda=\begin{bmatrix}3&0\\0&1\end{bmatrix}.
\]

核验：$\bm Q^\top\bm Q=\frac12\left[\begin{smallmatrix}2&0\\0&2\end{smallmatrix}\right]=\bm I$；又
\[
 \bm S\bm Q=\frac1{\sqrt2}\begin{bmatrix}3&1\\3&-1\end{bmatrix}=\bm Q\bm\Lambda ,
\]
第一列是$3\bm q_1$，第二列是$\bm q_2$，与$\bm\Lambda$的对角顺序一致。

结论：$\bm S=\bm Q\bm\Lambda\bm Q^\top$，即$\bm S$把$(1,1)^\top$方向放大$3$倍，保持$(1,-1)^\top$方向不变。
\end{example}

这个分解把$\bm S$的作用拆成三步：$\bm Q^\top$把向量转到特征坐标，$\bm\Lambda$在两个坐标方向上分别伸缩$3$倍和$1$倍，$\bm Q$再转回来。$\bm Q$与$\bm Q^\top$保持长度与夹角，改变形状的只有$\bm\Lambda$。由于两个方向的伸缩倍数不同，整个作用$\bm S$一般会改变向量间的夹角。例如$\bm e_1$与$\bm e_2$互相垂直，而$\bm S\bm e_1=(2,1)^\top$与$\bm S\bm e_2=(1,2)^\top$的内积为$4$，夹角余弦为$4/5$。只有特征方向之间的直角得以保持：$\bm S\bm q_1=3\bm q_1$与$\bm S\bm q_2=\bm q_2$仍然垂直。

上述路线适合手算小矩阵。它要求写出特征多项式并求出全部根，阶数较高时这一步本身就难以手工完成；大规模矩阵的数值特征值算法不在本节讨论。无论$\bm Q$与$\bm\Lambda$由哪种方法得到，上面两项核验都只用到矩阵乘法，同样适用：两项都成立，就足以推出$\bm S=\bm Q\bm\Lambda\bm Q^\top$。

\subsection{特征子空间与谱投影}
\label{sec:ls-spectral-projections}

特征值有重复时，$\bm Q$不唯一：重根空间内的任意一组标准正交基都可以作为这部分的列。本小节要回答，在这种不唯一之下，分解中哪些部分由$\bm S$本身确定。为此先把$\bm Q\bm\Lambda\bm Q^\top$按列展开，写成秩一矩阵之和：
\[
 \bm S=\sum_{i=1}^d\lambda_i\bm q_i\bm q_i^\top .
\]
由式~\eqref{eq:linear-algebra-orthogonal-projector}，$\bm q_i\bm q_i^\top$是到$\bm q_i$方向的正交投影。例如示例~\ref{ex:ls-symmetric-spectral-two-by-two}中
\[
 \bm S=\frac32\begin{bmatrix}1&1\\1&1\end{bmatrix}+\frac12\begin{bmatrix}1&-1\\-1&1\end{bmatrix}.
\]
单个$\bm q_i$在重根空间中可以改变，所以单项$\bm q_i\bm q_i^\top$也会随之改变。把同一特征值的各项合并，可以得到不再依赖这种选择的分解。

设$\bm S$的不同特征值为$\mu_1,\ldots,\mu_m$（$m\leq d$），特征空间为$E_j=\ker(\bm S-\mu_j\bm I)$。把$\bm Q$中满足$\lambda_i=\mu_j$的列排成矩阵$\bm Q_j$，令$\bm P_j=\bm Q_j\bm Q_j^\top$。

\begin{theorem}[谱投影]
\label{thm:ls-spectral-projections}
设$\bm S\in\mathbb R^{d\times d}$对称，记号如上。则：
\begin{enumerate}
 \item $E_j$由$\bm Q_j$的列张成，各$E_j$两两正交，且$\sum_j\dim E_j=d$；
 \item $\bm P_j$是到$E_j$的正交投影，满足$\bm P_j\bm P_k=\bm0$（$j\neq k$）与$\sum_{j=1}^m\bm P_j=\bm I$；
 \item $\bm S=\sum_{j=1}^m\mu_j\bm P_j$；
 \item $E_j$与$\bm P_j$只由$\bm S$决定，与重根空间内标准正交基的选取无关。
\end{enumerate}
\end{theorem}

\begin{proof}
(1)把任意$\bm x$写成$\bm x=\sum_ic_i\bm q_i$，则
\[
 (\bm S-\mu_j\bm I)\bm x=\sum_{i=1}^d c_i(\lambda_i-\mu_j)\bm q_i .
\]
$\bm q_i$线性无关，所以上式为零当且仅当$\lambda_i\neq\mu_j$时$c_i=0$，即$\bm x$属于$\bm Q_j$各列的张成。于是$E_j=\operatorname{range}\bm Q_j$。不同的$\bm Q_j$由$\bm Q$的不同列组成，所以$E_j$两两正交；各$\bm Q_j$的列数之和是$d$，所以维数之和为$d$。

(2)$\bm Q_j$的列是$E_j$的标准正交基，由式~\eqref{eq:linear-algebra-orthogonal-projector}，$\bm P_j$是到$E_j$的正交投影。$j\neq k$时$\bm Q_j^\top\bm Q_k=\bm0$，故$\bm P_j\bm P_k=\bm Q_j(\bm Q_j^\top\bm Q_k)\bm Q_k^\top=\bm0$。把$\bm Q\bm Q^\top=\sum_i\bm q_i\bm q_i^\top$的各项按特征值分组，得$\sum_j\bm P_j=\bm Q\bm Q^\top=\bm I$。

(3)在$\bm S=\sum_i\lambda_i\bm q_i\bm q_i^\top$中，把$\lambda_i=\mu_j$的项合并为$\mu_j\bm P_j$即可。

(4)$E_j=\ker(\bm S-\mu_j\bm I)$的定义只用到$\bm S$。设$\bm Q_j'$的列是$E_j$的另一组标准正交基，则每列都能用$\bm Q_j$的列表示，即$\bm Q_j'=\bm Q_j\bm R$，其中$\bm R=\bm Q_j^\top\bm Q_j'$。由$\bm R^\top\bm R=\bm Q_j'^\top\bm Q_j\bm Q_j^\top\bm Q_j'=\bm Q_j'^\top\bm P_j\bm Q_j'=\bm Q_j'^\top\bm Q_j'=\bm I$知$\bm R$是正交方阵，于是$\bm R\bm R^\top=\bm I$，
\[
 \bm Q_j'\bm Q_j'^\top=\bm Q_j\bm R\bm R^\top\bm Q_j^\top=\bm Q_j\bm Q_j^\top=\bm P_j .
\]
这里$\bm P_j\bm Q_j'=\bm Q_j'$是因为$\bm Q_j'$的列已在$E_j$中，投影不改变它们。
\end{proof}

称$\bm P_j$为$\bm S$的\term{谱投影}（spectral projection）。定理说明，重根空间内基的选择只是同一个投影的不同写法。下面的例子在三阶矩阵上直接看到这一点，同时演示重根空间内部的正交化。

\begin{example}[重根空间的两组基]
\label{ex:ls-repeated-eigenspace-projection}
给定
\[
 \bm S=\begin{bmatrix}2&1&1\\1&2&1\\1&1&2\end{bmatrix}=\bm I+\bm 1\bm 1^\top ,
\]
其中$\bm 1=(1,1,1)^\top$。

由$\bm S\bm 1=\bm 1+3\bm 1=4\bm 1$，$\bm 1$是特征值$4$的特征向量；对$\bm x\perp\bm 1$，$\bm S\bm x=\bm x$。所以特征值为$4$（单根）与$1$（二重），$E_1$是平面$x_1+x_2+x_3=0$。该平面的一组基$(1,-1,0)^\top$、$(1,0,-1)^\top$不正交，内积为$1$。Gram--Schmidt正交化后得
\[
 \bm q_2=\frac1{\sqrt2}\begin{bmatrix}1\\-1\\0\end{bmatrix},\qquad
 \bm q_3=\frac1{\sqrt6}\begin{bmatrix}1\\1\\-2\end{bmatrix}.
\]
直接相加，
\[
 \bm q_2\bm q_2^\top+\bm q_3\bm q_3^\top
 =\frac13\begin{bmatrix}2&-1&-1\\-1&2&-1\\-1&-1&2\end{bmatrix}
 =\bm I-\frac13\bm 1\bm 1^\top .
\]
若改从$(1,0,-1)^\top$出发正交化，另一组标准正交基为$(1,0,-1)^\top/\sqrt2$与$(1,-2,1)^\top/\sqrt6$；同样计算得到的仍是$\bm I-\frac13\bm 1\bm 1^\top$。

结论：$\bm P_4=\frac13\bm 1\bm 1^\top$，$\bm P_1=\bm I-\frac13\bm 1\bm 1^\top$，$\bm S=4\bm P_4+\bm P_1$。$E_1$中的基可以换，$\bm P_1$不变。
\end{example}

谱分解还能直接回答一个已经出现过的问题：给定对称矩阵，怎样判断它是否半正定或正定？第~\ref{sec:ls-positive-definite}节的定义要求检查所有$\bm x$上的二次型。把$\bm x$写成$\bm x=\sum_ic_i\bm q_i$，由正交性，式~\eqref{eq:matrix-analysis-quadratic-spectral-form}把二次型化为$\sum_i\lambda_ic_i^2$，交叉项全部消失，检查于是只剩下各特征值的符号。

\begin{theorem}[用特征值判定定号性]
\label{thm:ls-symmetric-definiteness-eigenvalues}
设$\bm S\in\mathbb R^{d\times d}$对称，特征值为$\lambda_1\geq\cdots\geq\lambda_d$。则$\bm S\succeq0$当且仅当$\lambda_d\geq0$，即全部特征值非负；$\bm S\succ0$当且仅当$\lambda_d>0$，即全部特征值为正。
\end{theorem}

\begin{proof}
取谱分解$\bm S=\bm Q\bm\Lambda\bm Q^\top$。必要性：取$\bm x=\bm q_i$，由式~\eqref{eq:matrix-analysis-quadratic-spectral-form}，$\bm q_i^\top\bm S\bm q_i=\lambda_i$。若$\bm S\succeq0$，则每个$\lambda_i\geq0$；若$\bm S\succ0$，由$\bm q_i\neq\bm0$得每个$\lambda_i>0$。

充分性：任取$\bm x=\sum_ic_i\bm q_i$，则$\sum_ic_i^2=\lVert\bm x\rVert_2^2$。若全部$\lambda_i\geq0$，则$\bm x^\top\bm S\bm x=\sum_i\lambda_ic_i^2\geq0$。若全部$\lambda_i>0$，则
\[
 \bm x^\top\bm S\bm x=\sum_{i=1}^d\lambda_ic_i^2\geq\lambda_d\sum_{i=1}^dc_i^2=\lambda_d\lVert\bm x\rVert_2^2 ,
\]
$\bm x\neq\bm0$时右端为正。
\end{proof}

第~\ref{sec:ls-svd}节构造奇异值时就要用到这个判据：$\bm A^\top\bm A$半正定，所以它的特征值都非负，可以开平方。证明中的$\bm q_i^\top\bm S\bm q_i=\lambda_i$只说明特征方向上的二次型值等于特征值；在全部单位方向中哪个方向使二次型最大或最小，是另一个问题，留到第~\ref{sec:ls-quadratic-extrema}节。

本节的结论可以归纳如下。能用正交矩阵对角化的实矩阵恰好是对称矩阵，谱分解$\bm S=\bm Q\bm\Lambda\bm Q^\top$把它写成沿互相垂直方向的伸缩，并可按“求特征值、求特征空间、空间内正交化、组装并核验”的次序算出。特征基在重根空间内可以改变，但特征子空间和谱投影$\bm P_j$由$\bm S$唯一确定，$\bm S=\sum_j\mu_j\bm P_j$；二次型在特征坐标下没有交叉项，定号性由特征值的符号决定。

许多重要的方阵并不对称，例如平面旋转与示例~\ref{ex:ls-triangular-power}的上三角矩阵$\left[\begin{smallmatrix}0.8&4\\0&0.5\end{smallmatrix}\right]$。下一节讨论，用正交矩阵换基时，一般方阵能化到多简单。
